\chapter{Generalization: Anti-Unification and MSG}
\label{chap:generalization}

\section{The Role of Generalization}
\label{sec:generalization:role}

When the whistle blows upon detecting that configuration $C_{\text{curr}}$ homeomorphically embeds an ancestor $C_{\text{anc}}$, the supercompiler must reconcile the two states. If the two states differ only by variable renaming, folding creates a direct recursion. However, if $C_{\text{curr}}$ contains growing sub-expressions (for instance, an accumulator variable growing from $0$ to $0 + 1$ to $0 + 1 + 2$), folding is impossible.

To prevent infinite unrolling while still discovering the underlying recursive invariant, the supercompiler must perform \emph{generalization}: finding a common, more general template that subsumes both configurations.

\section{The S\o{}rensen--Gl\"uck Most Specific Generalization Algorithm}
\label{sec:generalization:msg_alg}

NumLang implements the canonical S\o{}rensen and Gl\"uck (1995) anti-unification algorithm \citep{sorensen1995algorithm} in \texttt{src/mir/supercompiler/generalize.rs}.

Algorithm~\ref{alg:anti_unification} presents the formal anti-unification procedure.

\begin{algorithm}[h!]
\caption{The Most Specific Generalization (Anti-Unification) Algorithm}
\label{alg:anti_unification}
\begin{algorithmic}[1]
\Require Terms $t_1, t_2 \in \symterm$, Memoization Map $M : (\symterm \times \symterm) \to \mathcal{V}$
\Ensure Generalized Term $g$, Substitutions $\theta_1, \theta_2$
\Procedure{AntiUnify}{$t_1, t_2, M$}
    \If{$(t_1, t_2) \in M$}
        \State \Return $M(t_1, t_2)$
    \EndIf
    \If{$t_1 = f(u_1, \dots, u_n) \land t_2 = f(v_1, \dots, v_n)$}
        \ForAll{$i \in \{1, \dots, n\}$}
            \State $w_i \gets$ \Call{AntiUnify}{$u_i, v_i, M$}
        \EndFor
        \State \Return $f(w_1, \dots, w_n)$
    \Else
        \State Let $\alpha$ be a fresh symbolic variable
        \State $M \gets M \cup \{(t_1, t_2) \mapsto \alpha\}$
        \State $\theta_1 \gets \theta_1 \cup \{\alpha \mapsto t_1\}; \quad \theta_2 \gets \theta_2 \cup \{\alpha \mapsto t_2\}$
        \State \Return $\alpha$
    \EndIf
\EndProcedure
\end{algorithmic}
\end{algorithm}

\section{Variable Widening in the Interval Lattice}
\label{sec:generalization:widening}

When generalizing numeric state variables, naive anti-unification abstracts concrete bounds into unbounded fresh variables, destroying all downstream bounds check eliminations (BCE). 

NumLang couples syntactic MSG with an interval widening operator $\nabla$ over the interval domain $[lo, hi]$:
\begin{equation}
[lo_1, hi_1] \nabla [lo_2, hi_2] = \left[
\begin{cases}
lo_1 & \text{if } lo_1 \le lo_2 \\
-\infty & \text{otherwise}
\end{cases}, \quad
\begin{cases}
hi_1 & \text{if } hi_1 \ge hi_2 \\
+\infty & \text{otherwise}
\end{cases}
\right]
\end{equation}
By preserving stable lower bounds (such as array indices known to satisfy $i \ge 0$), NumLang retains bounds check elimination guarantees across generalization knots.

\section{Knot Materialization and Split vs. Generalize Decision}
\label{sec:generalization:split}

When an expression $e = f(t_1, t_2)$ triggers a whistle, the supercompiler faces a critical architectural decision:
\begin{enumerate}
    \item \textbf{Generalize}: Abstract growing subterms into variables and synthesize a generalized knot.
    \item \textbf{Split (Decompose)}: Decompose the compound expression into independent sub-expressions $t_1$ and $t_2$, driving each separately.
\end{enumerate}

NumLang implements Phase 48 \emph{Structural Whistle Decoupling}: the decision to split versus generalize is determined strictly by term size ratios and variable dependency metrics, never by string-based function name heuristics. If subterms $t_1$ and $t_2$ share zero free variables, splitting is chosen unconditionally, preventing unnecessary knot creation.

\section{The Anti-Unification Engine Implementation}
\label{sec:generalize:impl_listing}

The generalization engine computes Most Specific Generalizations (MSG) and widens states at knot points. Listing~\ref{lst:generalize_impl} reproduces the implementation from \texttt{src/mir/supercompiler/generalize.rs}:

\begin{lstlisting}[language=Rust, caption={Anti-Unification and Knot Materialization (\texttt{src/mir/supercompiler/generalize.rs}).}, label={lst:generalize_impl}]

use super::term::{SymTerm, SymTermId, TermInterner};
use crate::ast::BinaryOp;
use crate::mir::Place;
use crate::typecheck::types::Type;

#[derive(Debug, Clone)]
pub struct GeneralizationResult {
    pub common_term: SymTermId,
    pub var_mappings: Vec<(Place, SymTermId, SymTermId)>, // (fresh_var, t1_val, t2_val)
}

/// Computes the Most Specific Generalization (MSG) between two terms.
pub fn most_specific_generalization(
    t1: SymTermId,
    t2: SymTermId,
    interner: &mut TermInterner,
    next_var_id: &mut usize,
) -> GeneralizationResult {
    let mut memo: HashMap<(SymTermId, SymTermId), SymTermId> = HashMap::new();
    let mut mappings = Vec::new();

    let common = msg_helper(t1, t2, interner, next_var_id, &mut memo, &mut mappings);
    GeneralizationResult {
        common_term: common,
        var_mappings: mappings,
    }
}

/// Textbook first-order anti-unification (Sørensen & Glück 1995; Plotkin 1970).
pub fn anti_unify(
    t1: SymTermId,
    t2: SymTermId,
    interner: &mut TermInterner,
    subst1: &mut HashMap<String, SymTermId>,
    subst2: &mut HashMap<String, SymTermId>,
) -> SymTermId {
    let mut next_var_id = subst1.len();
    let res = most_specific_generalization(t1, t2, interner, &mut next_var_id);
    for (place, v1, v2) in res.var_mappings {
        subst1.insert(place.local.clone(), v1);
        subst2.insert(place.local, v2);
    }
    res.common_term
}

fn msg_helper(
    t1: SymTermId,
    t2: SymTermId,
    interner: &mut TermInterner,
    next_var_id: &mut usize,
    memo: &mut HashMap<(SymTermId, SymTermId), SymTermId>,
    mappings: &mut Vec<(Place, SymTermId, SymTermId)>,
) -> SymTermId {
    if t1 == t2 {
        return t1;
    }

    if let Some(&cached) = memo.get(&(t1, t2)) {
        return cached;
    }

    let term1 = interner.get(t1).clone();
    let term2 = interner.get(t2).clone();

    let result = match (&term1, &term2) {
        (SymTerm::Binary(op1, l1, r1, ty1), SymTerm::Binary(op2, l2, r2, _)) if op1 == op2 => {
            let common_l = msg_helper(*l1, *l2, interner, next_var_id, memo, mappings);
            let common_r = msg_helper(*r1, *r2, interner, next_var_id, memo, mappings);
            interner.intern_binary(*op1, common_l, common_r, ty1.clone())
        }
        (SymTerm::Unary(op1, in1, ty1), SymTerm::Unary(op2, in2, _)) if op1 == op2 => {
            let common_in = msg_helper(*in1, *in2, interner, next_var_id, memo, mappings);
            interner.intern_unary(*op1, common_in, ty1.clone())
        }
        (SymTerm::Constructor(n1, t1, f1, ty1), SymTerm::Constructor(n2, t2, f2, _))
            if n1 == n2 && t1 == t2 && f1.len() == f2.len() =>
        {
            let mut common_f = Vec::with_capacity(f1.len());
            for (&a, &b) in f1.iter().zip(f2.iter()) {
                common_f.push(msg_helper(a, b, interner, next_var_id, memo, mappings));
            }
            interner.intern_constructor(n1.clone(), *t1, common_f, ty1.clone())
        }
        (SymTerm::Call(c1, a1, ty1), SymTerm::Call(c2, a2, _))
            if c1 == c2 && a1.len() == a2.len() =>
        {
            let mut common_a = Vec::with_capacity(a1.len());
            for (&a, &b) in a1.iter().zip(a2.iter()) {
                common_a.push(msg_helper(a, b, interner, next_var_id, memo, mappings));
            }
            interner.intern_call(c1.clone(), common_a, ty1.clone())
        }
        (SymTerm::ClosureVal(fn1, c1, ty1), SymTerm::ClosureVal(fn2, c2, _))
            if fn1 == fn2 && c1.len() == c2.len() =>
        {
            let mut common_c = Vec::with_capacity(c1.len());
            for (&a, &b) in c1.iter().zip(c2.iter()) {
                common_c.push(msg_helper(a, b, interner, next_var_id, memo, mappings));
            }
            interner.intern_closure_val(fn1.clone(), common_c, ty1.clone())
        }
        (SymTerm::Thunk(b1, c1, ty1), SymTerm::Thunk(b2, c2, _))
            if b1 == b2 && c1.len() == c2.len() =>
        {
            let mut common_c = Vec::with_capacity(c1.len());
            for (&a, &b) in c1.iter().zip(c2.iter()) {
                common_c.push(msg_helper(a, b, interner, next_var_id, memo, mappings));
            }
            interner.intern_thunk(b1.clone(), common_c, ty1.clone())
        }
        (SymTerm::Ref(i1, ty1), SymTerm::Ref(i2, _)) => {
            let common = msg_helper(*i1, *i2, interner, next_var_id, memo, mappings);
            interner.intern_ref(common, ty1.clone())
        }
        (SymTerm::Deref(p1, ty1), SymTerm::Deref(p2, _)) => {
            let common = msg_helper(*p1, *p2, interner, next_var_id, memo, mappings);
            interner.intern_deref(common, ty1.clone())
        }
        (SymTerm::Discriminant(i1, _), SymTerm::Discriminant(i2, _)) => {
            let common = msg_helper(*i1, *i2, interner, next_var_id, memo, mappings);
            interner.intern_discriminant(common)
        }
        (SymTerm::Select(c1, t1, e1, ty1), SymTerm::Select(c2, t2, e2, _)) => {
            let common_c = msg_helper(*c1, *c2, interner, next_var_id, memo, mappings);
            let common_t = msg_helper(*t1, *t2, interner, next_var_id, memo, mappings);
            let common_e = msg_helper(*e1, *e2, interner, next_var_id, memo, mappings);
            interner.intern_select(common_c, common_t, common_e, ty1.clone())
        }
        _ => {
            // Generalize to fresh variable
            let var_name = format!("_gen_{}", *next_var_id);
            *next_var_id += 1;
            let ty = term1.ty().clone();
            let place = Place {
                local: var_name,
                projections: vec![],
            };
            let var_id = interner.intern_var(place.clone(), ty);
            mappings.push((place, t1, t2));
            var_id
        }
    };

    memo.insert((t1, t2), result);
    result
}

/// Solves closed forms for numerical recurrence sequences:
/// Given simulated initial states [s0, s1, s2, s3], detects polynomial degree
/// and returns closed form in terms of `num_iters`.
pub fn solve_recurrence(
    samples: &[i64],
    num_iters: SymTermId,
    interner: &mut TermInterner,
) -> Option<SymTermId> {
    if samples.len() < 3 {
        return None;
    }

    let s0 = samples[0];
    let s1 = samples[1];
    let s2 = samples[2];

    // First differences
    let d1_0 = s1.wrapping_sub(s0);
    let d1_1 = s2.wrapping_sub(s1);

    // Degree 1: Linear sequence s_k = s0 + k * d1
    if d1_0 == d1_1 {
        // s0 + N * d1
        let s0_term = interner.intern_int(s0);
        let d1_term = interner.intern_int(d1_0);
        let n_times_d1 = interner.intern_binary(
            BinaryOp::Mul,
            num_iters,
            d1_term,
            Type::I64,
        );
        return Some(interner.intern_binary(
            BinaryOp::Add,
            s0_term,
            n_times_d1,
            Type::I64,
        ));
    }

    // Degree 2: Quadratic / Triangular sequence
    if samples.len() >= 4 {
        let s3 = samples[3];
        let d1_2 = s3.wrapping_sub(s2);
        let d2_0 = d1_1.wrapping_sub(d1_0);
        let d2_1 = d1_2.wrapping_sub(d1_1);

        if d2_0 == d2_1 {
            // Formula: s_k = s0 + k * d1_0 + (k * (k - 1) / 2) * d2_0
            let s0_term = interner.intern_int(s0);
            let d1_term = interner.intern_int(d1_0);
            let d2_term = interner.intern_int(d2_0);
            let one_term = interner.intern_int(1);
            let two_term = interner.intern_int(2);

            // k * d1_0
            let linear_part = interner.intern_binary(
                BinaryOp::Mul,
                num_iters,
                d1_term,
                Type::I64,
            );

            // (k - 1)
            let k_minus_1 = interner.intern_binary(
                BinaryOp::Sub,
                num_iters,
                one_term,
                Type::I64,
            );

            // k * (k - 1)
            let k_times_k_minus_1 = interner.intern_binary(
                BinaryOp::Mul,
                num_iters,
                k_minus_1,
                Type::I64,
            );

            // k * (k - 1) / 2
            let tri_part = interner.intern_binary(
                BinaryOp::Div,
                k_times_k_minus_1,
                two_term,
                Type::I64,
            );

            // (k * (k - 1) / 2) * d2_0
            let quad_part = interner.intern_binary(
                BinaryOp::Mul,
                tri_part,
                d2_term,
                Type::I64,
            );

            // s0 + linear + quad
            let part1 = interner.intern_binary(
                BinaryOp::Add,
                s0_term,
                linear_part,
                Type::I64,
            );
            return Some(interner.intern_binary(
                BinaryOp::Add,
                part1,
                quad_part,
                Type::I64,
            ));
        }
    }

    // Degree 3: Cubic sequence
    if samples.len() >= 5 {

\end{lstlisting}
